\subsection{Every single pair of minimum three: the reference proof}\label{proof:R2-C032}
\paragraph{R32.1 The original object and directions.}
All operators, products and expectations use the original finite probability
measure \(\mu\) and a real skew kernel with \(|K(x,\allowbreak y)|\le1\). Directions
are chosen from the original tree before the host. Original arm
multiplicities, subdivisions and the four independent positive core lengths
are retained.

\emph{Theorem A.} For every receiving arm \(a\ge3\), terminal residual
list exactly \((3,\allowbreak b)\), \(b\ge4\), and original connector length \(L\ge1\), use
\fitdisplay{
\begin{array}{c|c|c}
L&\text{Actual terminal }F&\text{All connector edges}\\\hline
1\text{ or }2&f_3\bar f_b&\text{terminal to receiver}\\
L\ge3&h_3f_b&\text{terminal to receiver}.
\end{array}
}{}
The receiver's existing arm points outward, giving \(f_a\). An \(f\)-arm
is fully outward, a \(\bar f\)-arm fully inward, and \(h_3\) has its first
edge toward the terminal and its other two edges toward the leaf. For
genuine equal-arm-pair products \(E_o,\allowbreak E_t\) at their original locations,
\fitdisplay{
\E\bigl[E_of_a\{T^L(E_tF)\}^4\bigr]\le1-d/3000.
}{}
The stronger deficits are \(d/10\) for \(L=1\), \(3d/320\) for \(L=2\),
and the explicit low/high deficits below for later lengths. The endpoint retains its two unequal arms.

\emph{Theorem B.} The component statements listed below complete the
directory for every strictly increasing positive terminal residual list
of size at least two, every original \(a\ge3,\allowbreak L\ge1\), with
common same-\(f_a\) deficit \(d/3000\). All short-length direction
exceptions in those statements are retained.

Single pairs of minimum one or two and unit-minimum odd triples use
Section~\ref{proof:R2-C009}; minimum-three single pairs use Theorem A.
Even lists of minimum at least four use Section~\ref{proof:R2-C022};
even lists with at least two pairs and minimum two or three use
Section~\ref{proof:R2-C025}; and even multipair lists of minimum one
use Section~\ref{proof:R2-C028}.
Odd lists with nonunit minimum use Section~\ref{proof:R2-C018} when
the second length is three and Section~\ref{proof:R2-C020} otherwise.
For unit-minimum odd lists with at least five entries, second length
two, three or at least four uses, respectively,
Sections~\ref{proof:R2-C024}, \ref{proof:R2-C023} and
\ref{proof:R2-C021}. These alternatives cover every strict residual
list, preserve the actual outward receiving arm \(f_a\), and choose
directions from the original integer data. Every cited deficit is at
least \(d/3000\).

The central allocations using these local reference inputs are
given in Section~\ref{sec:five-star-fork}.

\paragraph{R32.2 Original finite calculus and receiver.}
Delete only zero-mass host atoms and put
\fitdisplay{
\begin{gathered}
Jv=(\E v)\one,\quad B=-K,\quad u=B\one,\quad T=J+B,\quad A=J-B=T^*,\\
p=1-u,\quad q=1+u,\quad f_n=A^n\one,\quad\bar f_n=T^n\one,\\
h_n=Tf_{n-1},\quad W_n=A\bar f_{n-1},\quad
\delta_i=\lVert B^i\one\rVert_2^2,\quad d=\delta_1,\\
\theta=13/32,\quad\kappa=19/32,\quad\tau=d-\delta_2.
\end{gathered}
}{}
The finite calculus of Appendix~\ref{q3:app:foundations} gives positive \(A,\allowbreak T\) and
\fitdisplay{
\begin{gathered}
0\le d\le1/3,\quad d^2\le\delta_2\le\theta d,\quad
\delta_{i+1}\le\theta\delta_i,\quad\lVert B\rVert^2\le\theta,\\
|Bv(x)|\le\lVert v\rVert_2,\qquad p(x),q(x)\ge\mu(x)>0.
\end{gathered}
}{}
Also \(d^2\le\delta_2\) follows directly from \(\E Bu=-d\) and
Cauchy--Schwarz. The finite sums are
\fitdisplay{
G_n(z)=\sum_{j=0}^{n-1}z^j,\quad G_0=0,\quad\beta_1=\beta_2=0,
\qquad
\beta_n=\sum_{j=1}^{\lfloor(n-1)/2\rfloor}(-1)^{j-1}\delta_j\quad(n\ge3).
}{}
The exact positive convolution and its full reverse are
\fitdisplay{
\begin{aligned}
1+G_n(B)u&=\bar f_n+\sum_{r=3}^n\beta_r\bar f_{n-r},\\
1-G_n(-B)u&=f_n+\sum_{r=3}^n\beta_rf_{n-r},
\qquad\tau\le\beta_r\le d\quad(r\ge3).
\end{aligned}
}{}
The sums are empty for \(n=1,\allowbreak 2\). The pure path messages \(f_n\) and \(\bar f_n\) lie in \([0,\allowbreak 2]\) and have means at most one. The identity \(f_\ell+h_\ell=2\E f_{\ell-1}\)
gives \(0\le f_\ell h_\ell\le1\); reversal gives
\(0\le\bar f_\ell W_\ell\le1\). These are the freely removable genuine
pair decorations.

The original spectral measure \(\nu\) of \(-B^2\) at \(u\) has mass \(d\),
support \([0,\allowbreak \theta]\), and \(\int x^j\,\allowbreak d\nu=\delta_{j+1}\).
Integrals below use this measure. At \(z=i\sqrt x\), for a real finite
polynomial \(Q=A_0+zC_0\),
\fitdisplay{
\begin{aligned}
\lVert Q(B)u\rVert_2^2&=\int(A_0^2+xC_0^2),\\
\E[Q(B)u]&=-\int C_0,\qquad
\langle Q(B)u,u\rangle=\int A_0.
\end{aligned}
}{}
Moments such as \(\tau\) remain global. The ratio \(\delta_2/d\) does not impose a pointwise condition \(x\ge d\) on the spectral support.
For \(G_n=A_n+zC_n\), every finite \(n\ge1\) has
\(1-x\le A_n\le1\), \(0\le C_n\le1\), with \((A_1,\allowbreak C_1)=(1,\allowbreak 0)\).
For \(n\ge2\), both coordinates are at least \(1-x\). For \(n\ge3\),
\(A_n\le1-x+x^2\), and for \(n\ge4\), the same holds for \(C_n\).
Finite parity gives \(C_n-A_n\le x\). The exact geometric formula gives
\(\lVert G_n\rVert\le1\) for \(n=1\) or \(n\ge3\), whereas
\(\lVert G_2\rVert^2\le1+\theta<(19/16)^2\). Every finite row therefore
has norm less than \(19/16\), with both \(G_1,\allowbreak G_2\) exceptions retained.
We also have the finite gaps
\fitdisplay{
|G_4|\le1-x/2,\qquad |G_5|\le1-\kappa x/2,\qquad
|G_n|\le1-7x/30\quad(n\ge6).
}{}
For \(G_4\), use \((1-x)^2(1+x)\); for \(G_5\), its exact square
\(1-x+x^2-x^3+x^4\le1-\kappa x\). More generally,
\(|G_n|^2\le1-7x/15\) for \(n\ge3\). The cases \(3,\allowbreak 4\) are direct;
for \(n\ge5\), the finite numerator is at most \((1+x^3)^2\).
The polynomial \(8-7x-30x^2-15x^5\) decreases on \([0,\allowbreak \theta]\) and
equals \(1311885/33554432>0\) at \(\theta\). Taking square roots gives
the gaps, with the tails bounding the finite sums.

For every \(a\ge3\), the true weak receiving upper is
\fitdisplay{
\begin{aligned}
C&=1-G_a(-B)u-\beta_a\ge f_a\ge0,\\
\E C&=1-\beta_a-\beta_{a+1}
 \le1-2\tau\le1-19d/16\le1-d,\\
C&=p-\beta_a+BG_{a-1}(-B)u\le p+(6/5)\sqrt d.
\end{aligned}
}{}
For \(H=Q(B)u\) of actual mean \(-\alpha\),
\fitdisplay{
\operatorname{Cov}(C,H)
=-\int\operatorname{Re}(G_aQ)-\beta_{a+1}\alpha.
}{}
Indeed \(\E[G_a(-B)u]=\beta_{a+1}\), fixing the sign before estimating.
The \(a=3\) predecessor is the paid \(G_2\) row. Uniformly in \(a\ge3\),
\fitdisplay{
\begin{array}{c|c}
n&\text{Lower bound for }\operatorname{Re}(G_aG_n)\\\hline
2&1-2x\ge3/16\\
3&1-3x+x^2\ge-55/1024>-3/50\\
4&(1-x)(1-2x)\ge57/512\\
5&(1-x)^3>1/5\\
n\ge6&1-3x+2x^2-x^3\ge1451/32768.
\end{array}
}{}
At \(n=6\), \(G_6=(1-x+x^2)(1+z)\) gives the stronger \(57/512\).
For \(n\ge6\), use \((1-x)^2-x(1-x+x^2)\); for \(n=5\), use
\((1-x)(1-x+x^2)-x(1-x)\). The other rows are immediate, and the
polynomials decrease on the interval. The general tail row is
\(1451/32768\), not \(1/10\). Upper real products are at most one when
used: for \(n=2\), \(A_a-xC_a\le1\); for \(n\ge3\), use the two
contractions.

\paragraph{R32.2.1 Fourth expansion and finite positive propagation.}
For nonnegative \(Y\) of mean \(m_Y\),
\fitdisplay{
\begin{aligned}
\E(CY^4)={}&(\E C)m_Y^4+4m_Y^3\operatorname{Cov}(C,Y)\\
&+\E\bigl[C(Y-m_Y)^2\{(Y+m_Y)^2+2m_Y^2\}\bigr].
\end{aligned}
}{}
For \(W=e^{-\Gamma}Y\), if the actual scaled curvature is paid by \(M\),
the variance contribution is \(Me^{-2\Gamma}\operatorname{Var}Y\), with
second-power scaling. For a canonical \(F\le e^{R-\Gamma}\), put \(D_i=\E B^iR\) and choose
\(D\ge D_0+(6/5)\E R^2\) when \(|R|\le1/2\). The factorial tail for
\(R\ge0\), and integral remainder for \(R\le0\), give
\(e^R\le1+R+(3/5)R^2\). Positivity and \(T(R^2)\le2\E R^2\) imply
\fitdisplay{
0\le e^\Gamma T^LF\le Y_L
=\bar f_L+B^LR+D\bar f_{L-1}
 +\sum_{i=1}^{L-1}D_i\bar f_{L-1-i}.
}{}
Write \(R=U(B)u\), \(\epsilon=D+D_1\), and, for \(n\ge2\),
\(P_n=G_n+z^nU\), \(Q_n=P_n+\epsilon G_{n-2}+Dz^{n-2}\).
Finite substitution gives
\fitdisplay{
\begin{aligned}
Y_L&=1+\epsilon+Q_L(B)u-\sum_{n=3}^L\gamma_n\bar f_{L-n},\\
\gamma_3&=d-D_2,\qquad\gamma_4=d+Dd-D_3,\\
\gamma_n&=\beta_n+D\beta_{n-1}+D_1\beta_{n-2}-D_{n-1}\quad(n\ge5),\\
\E[Q_n(B)u]&=-\gamma_{n+1},\qquad
\langle Q_n(B)u,u\rangle=\gamma_{n+2}.
\end{aligned}
}{}
The last identity follows by taking expectations in
\(Q_{n+1}(B)u=(1+\epsilon)u+BQ_n(B)u\) and using skewness; its future
index differs from the output mean. In particular,
\(Y_3=(1+\epsilon)q-\gamma_3+BQ_2(B)u\).
If \(\gamma_n\ge hd\) for all \(n\ge3\), retain the last two genuine
convolution terms for every \(L\ge4\) to obtain the positive upper
\fitdisplay{
\begin{aligned}
Z_L&=1+\epsilon+Q_L(B)u-hd(1+q)\\
&=(1+\epsilon-hd)q-hd+BQ_{L-1}(B)u,\\
\E Z_L&\le1+\epsilon-3hd,\\
\operatorname{Var}Z_L
&=\lVert Q_L(B)u\rVert_2^2-2hd\gamma_{L+2}+h^2d^3-\gamma_{L+1}^2\\
&\le\lVert Q_L\rVert^2d.
\end{aligned}
}{}
All other removed terms have nonnegative original path factors and proved
nonnegative coefficients. Both future indices remain paid. 
\paragraph{R32.3 Global \(L=1\), with actual terminal \(f_3\bar f_b\).}
Fix \(a\ge3,\allowbreak b\ge4\). The short arm points outward, the long arm inward,
the one connector terminal-to-receiver, and the receiver is \(f_a\).
We prove \(\E[f_a\{T(f_3\bar f_b)\}^4]\le1-d/10\) on \(0\le d\le1/3\).
Use the original receiver \(C\), measure \(\nu\) and covariance formula
above throughout.

\paragraph{R32.A1 Actual positive square and its negative constant.}
Put \(n=b-3\ge1\), \(G_n=A_0+zC_0\). For every finite \(n\),
\(1-x\le A_0\le1\), \(0\le C_0\le1\); at \(n=1\), \((A_0,\allowbreak C_0)=(1,\allowbreak 0)\).
The exact arm and positive convolution give
\fitdisplay{
f_3=1-d-u+Bu-B^2u,\qquad
\bar f_b\le1+G_b(B)u-\beta_b.
}{}
Consequently
\fitdisplay{
\begin{aligned}
f_3+\bar f_b&\le2-\Gamma+R,\qquad
\Gamma=d+\beta_b\in[2d-\delta_2,2d],\\
R&=2Bu+B^3G_n(B)u=[x^2C_0+zr]u,\\
r&=2-xA_0\in[2-x,2-x+x^2].
\end{aligned}
}{}
Bracket notation denotes the finite spectral polynomial applied to \(u\).
The nonnegative actual arm sum gives \(m+R/2\ge0\), where
\(m=1-\Gamma/2\in[2/3,\allowbreak 1]\), \(k=1/m\in[1,\allowbreak 3/2]\). Hence
\fitdisplay{
F\le m^2(1+kR+k^2R^2/4),\qquad
0\le TF\le m^2Y\le e^{-\Gamma}Y,
}{}
where \(Y=q+kBR+\delta\), \(\delta=k\E R+(k^2/2)\E R^2\).
This uses \(T(R^2)\le2\E R^2\), with no exponential Taylor approximation
of \(R\). The actual means are
\fitdisplay{
\E R=-\int r\le-51d/32,\qquad \E BR=-\int x^2C_0\le0.
}{}
The full square \(\sigma=\E R^2\) satisfies
\fitdisplay{
\sigma\le\int P(x),\qquad P=x(2-x+x^2)^2+x^4,
}{}
and \(P'=(2-x+x^2)(2-3x+5x^2)+4x^3>0\). Its value at \(\theta\) is
\(43080765/33554432\), smaller than \(21/16\) by \(959427/33554432\).
Thus \(\sigma<21d/16\) and
\fitdisplay{
\delta/k\le-51d/32+(3/4)(21/16)d=-39d/64.
}{}
In particular \(\delta\le-3d/5\) and
\(0\le m_Y=\E Y=1+\delta+k\E BR\le1-3d/5\).
Pointwise,
\(Y\le q+k\sqrt\sigma\le q+(7/4)\sqrt d\), since \(k\le3/2\) and
\(3\sqrt{21}<14\). It is nonnegative because it dominates \(TF/m^2\).

\paragraph{R32.A2 Complete variance contraction for \(1\le k\le3/2\).}
The random part of \(Y\) is \(Q_k(B)u\), with
\(Q_k=(1-kxr)+zkx^2C_0\). Its squared modulus is convex in \(k\)
at fixed \(x,\allowbreak A_0,\allowbreak C_0\), so consider \(k=1,\allowbreak 3/2\).
At \(k=1\), the real coordinate is positive and
\fitdisplay{
|Q_1|^2\le(1-x)^4+x^5\le1-2x.
}{}
The last difference divided by \(x\) is \(-2+6x-4x^2+x^3+x^4\).
It increases on \([0,\allowbreak 5/12]\) to \(-1907/20736<0\); \(x=0\) is direct.
At \(k=3/2\), the squared difference from \(1-2x\), divided by \(x\),
is at most \(2-3r+(9/4)xr^2+(9/4)x^4\), convex in \(r\).
At its two allowed endpoints it equals
\fitdisplay{
\begin{aligned}
J_{\rm lo}&=-4+12x-9x^2+(9/4)x^3+(9/4)x^4,\\
J_{\rm hi}&=-4+12x-12x^2+(45/4)x^3-(9/4)x^4+(9/4)x^5.
\end{aligned}
}{}
Both increase on \([0,\allowbreak 5/12]\), with derivatives at least \(9/2\) and
\(259/192\), respectively. Their endpoint values are \(-3059/9216\)
and \(-34183/110592\), both negative. This pays every finite \(n\)
and every \(k\), so, retaining the actual mean before weakening,
\fitdisplay{
\operatorname{Var}Y=\lVert Q_k(B)u\rVert_2^2-(k\E BR)^2\le d-2\delta_2.
}{}
Writing \(G_a=A_a+zC_a\), its coordinates obey \(1-x\le A_a\le1\) and
\(0\le C_a\le1\). Without assuming the real \(Q_k\)
coordinate positive,
\fitdisplay{
\begin{aligned}
\operatorname{Re}(G_aQ_k)
&\ge1-x-kx(2-x+x^2)-kx^3\\
&=1-(1+2k)x+kx^2(1-2x)\ge1-4x.
\end{aligned}
}{}
Its mean is \(k\E BR\le0\), so the receiver mean-product is favorable.
Hence \(\operatorname{Cov}(C,\allowbreak Y)\le-d+4\delta_2\); this possibly positive
upper is retained.

\paragraph{R32.A3 The scaled coefficient \(10\).}
Put \(t=\sqrt d\le3/5\), \(W=e^{-\Gamma}Y\). As \(\E Y\le1\) and
\(Y\le q+7t/4\), the same-reference remainder coefficient is at most
\(e^{-2\Gamma}H\), where
\fitdisplay{
H=(2-q+6t/5)\{(q+1+7t/4)^2+2\}.
}{}
Its constant is \(6+q-q^3\le32/5\), and
\fitdisplay{
\begin{aligned}
H_1&=53/5+(59/10)q-(23/10)q^2\le29/2,\\
H_2&=413/40+(91/80)q\le63/5,\qquad
H_3=147/40,\quad H_4=H_5=0.
\end{aligned}
}{}
The first maximum has margin \(107/920\) below \(29/2\). Hence
\fitdisplay{
\begin{aligned}
H&\le32/5+(29/2)t+(2961/200)t^2,\\
10e^{2\Gamma}-H&\ge18/5-(29/2)t+(1707/100)t^2>0.
\end{aligned}
}{}
Here \(\Gamma/d\ge51/32\); the last discriminant is \(-17779/500\).
This proves the continuous bound, and the exact fourth expansion gives
\fitdisplay{
\begin{aligned}
\E(CW^4)\le{}&e^{-4\Gamma}
 [(1-d)m_Y^4+4m_Y^3\operatorname{Cov}(C,Y)]\\
&+10e^{-2\Gamma}\operatorname{Var}Y.
\end{aligned}
}{}
The variance term therefore has the second-power scale.

\paragraph{R32.A4 Coupled scalar payment on the original domain.}
For \(d>0\), let \(\bar x=\delta_2/d\), so \(d\le\bar x\le\theta\).
The true mean upper \(m_Y\le1-3d/5\), covariance and variance bounds,
and \(\Gamma\ge2d-\delta_2\), bound the moment by
\fitdisplay{
\begin{aligned}
F(d,\bar x)={}&e^{-(57/5-4\bar x)d}
 +10(1-2\bar x)d e^{-(4-2\bar x)d}\\
&+4(4\bar x-1)_+d e^{-(49/5-4\bar x)d}.
\end{aligned}
}{}
The positive part bounds the covariance with either sign. For fixed \(0<d\le1/3\), \(F\) decreases in \(\bar x\) on
\([0,\allowbreak \theta]\). Below \(1/4\), its derivative divided by
\(d e^{-(4-2\bar x)d}\) is
\(4e^{-(37/5-2\bar x)d}-20+20d(1-2\bar x)<0\).
Above \(1/4\), add \(16e^{-yd}(1+zd)\), where
\(y=29/5-2\bar x\), \(z=4\bar x-1\). We have \(4<y<6\),
\(0\le z<1\), \(y-z>3\), and
\fitdisplay{
e^{-yd}(1+zd)\le\frac{1+zd}{1+yd}
\le1-\frac{3d}{1+6d}\le1-d.
}{}
The full bracket is at most \(4-20+10d+16(1-d)=-6d<0\).
Continuity at \(1/4\) gives \(F(d,\allowbreak \bar x)\le F(d,\allowbreak d)\), using \(\delta_2\ge d^2\).
For \(0<d\le1/4\), the positive part vanishes at \(\bar x=d\).
Set \(A=57/5-4d\), \(B=4-2d\), \(C_0=A-B=37/5-2d\).
Positive quadratic Taylor for \(e^{Bd}\) and the quadratic upper for
\(e^{-C_0d}\) give
\fitdisplay{
\frac{e^{Bd}-e^{-C_0d}}d-10(1-2d)
\ge7/5-(169/50)d+(34/5)d^2\ge111/200.
}{}
Since \(Bd\le1\), \(e<3\), this gives \(F(d,\allowbreak d)\le1-37d/200\).
For \(1/4\le d\le1/3\), bound its three terms separately by
\fitdisplay{
\begin{gathered}
\text{first}<e^{-5/2}<1/8,\\
\text{second}\le5d e^{-10d/3}\le3/(2e)<3/5,\\
\text{third}<4/(9e^2)<4/63<1/15.
\end{gathered}
}{}
The inequalities \(e>5/2\), \(e^2>7\) follow from positive Taylor
series. The sum is less than \(19/24\le1-5d/8\). At \(d=0\), every
actual message is one even if \(K\ne0\). Thus
\fitdisplay{
\E[f_a\{T(f_3\bar f_b)\}^4]\le1-d/10
\qquad(a\ge3,\ b\ge4).
}{}
Genuine decorations restore by positivity and full reversal reverses
every original edge together.

\paragraph{R32.4 Global \(L=2\), with the same \(f_3\bar f_b\).}
Use the actual terminal above. Put
\(\Gamma=d+\beta_b\), \(m=1-\Gamma/2\), \(k=1/m\), \(a_0=1-3d/5\).
Then \(\Gamma\ge51d/32\), \(2/3\le m\le1\), \(1\le k\le3/2\),
\(4/5\le a_0\le1\). The exponent-free envelope is
\fitdisplay{
\begin{aligned}
F&\le m^2(1+kR+k^2R^2/4),\\
R&=2Bu+B^3G_{b-3}(B)u=[x^2C_0+zr]u,\\
r&=2-xA_0,\quad1-x\le A_0\le1,\quad0\le C_0\le1,\\
2-x&\le r\le2-x+x^2\le2.
\end{aligned}
}{}
The signed means are \(\E R=-\int r\),
\(D_1=\E BR=-\int x^2C_0\le0\), \(D_2=\E B^2R=\int xr\).
The same-resource bounds give \(\E R\le-51d/32\), \(\sigma\le21d/16\)
and \(\delta=k\E R+(k^2/2)\sigma\le-39d/64\le-3d/5\).

\paragraph{R32.E1 Positive propagation and complete mean.}
Positivity first gives \(TF\le m^2Y_1\),
\(Y_1=q+kBR-3d/5\ge0\). Apply the second actual \(T\):
\fitdisplay{
\begin{aligned}
T^2F&\le m^2Y,\qquad Y=q+\epsilon+BV\ge0,\\
\epsilon&=-3d/5+kD_1,\qquad V=-(3d/5)\one+u+kBR,\\
Q_2&=(a_0-kx^3C_0)+z(1-kxr),\qquad Y=1+\epsilon+Q_2(B)u.
\end{aligned}
}{}
It is positive as the actual \(T\)-image of \(Y_1\); no signed upper is
squared. Its mean is
\(\mu_Y=1-8d/5+k\int(xr-x^2C_0)\). The bracket is nonnegative because
\(r-xC_0\ge2-2x>0\), while
\(xr\le x(2-x+x^2)\le\theta(2-\theta+\theta^2)=23413/32768\).
Therefore
\fitdisplay{
1-8d/5\le\mu_Y
\le1-[8/5-70239/65536]d\le1-21d/40.
}{}
The last difference is \(1061/327680>0\), and \(\mu_Y\ge7/15\).

\paragraph{R32.E2 Precursor, variance and covariance.}
The actual precursor contains the constant vector \(-(3d/5)\one\).
Let \(c=3d/5\). R32.A2 gives \(\lVert u+kBR\rVert_2^2\le d-2\delta_2\).
Since \(\E(u+kBR)=kD_1\), \(D_1=-\int x^2C_0\), its complete square is
\fitdisplay{
\begin{aligned}
\lVert V\rVert_2^2
&=\lVert u+kBR\rVert_2^2+c^2-2ckD_1\\
&\le d-2\delta_2+9d^2/25+(9/5)d\delta_3\\
&\le d-(281/160)\delta_2+9d^2/25\\
&\le d-(1117/800)d^2\le d.
\end{aligned}
}{}
This uses \(\delta_3\le\theta\delta_2\), \(d\le1/3\), \(\delta_2\ge d^2\),
and retains the positive cross \(-2ckD_1\). The original row now gives
\(Y\le q+\epsilon+\sqrt d\le q-3d/5+\sqrt d\le q+\sqrt d\).
The output polynomial \(Q_2\) is used separately for variance/covariance;
the precursor remains the actual constant vector, rather than a multiple
of \(u\).

The real coordinate of \(Q_2\) is positive and at most \(a_0\), since
\(a_0-kx^3C_0\ge4/5-(3/2)\theta^3>0\). Put \(y=kxr\).
For \(y\le1\),
\(x(1-y)^2\le x(1-x)^4\le256/3125<1/12\); the derivative is
\((1-x)^3(1-5x)\), with maximum at \(x=1/5\).
For \(y\ge1\), \(y-1\le4703/65536<3/40\), so
\(x(1-y)^2<9/3200<1/12\). Before the favorable squared-mean subtraction,
\fitdisplay{
\operatorname{Var}Y\le\lVert Q_2(B)u\rVert_2^2
\le(a_0^2+1/12)d\le[13/12-27d/25]d,
}{}
using \(9d^2/25\le3d/25\).
For the actual receiver,
\fitdisplay{
\operatorname{Cov}(C,Y)
=-\int\operatorname{Re}(G_aQ_2)-\beta_{a+1}\alpha,
\qquad\alpha=d-kD_2.
}{}
For every finite \(a\ge3\), \(\operatorname{Re}(G_aQ_2)\ge7/25\).
If \(1-kxr\ge0\), the nonnegative coordinates give
\fitdisplay{
\begin{aligned}
\operatorname{Re}(G_aQ_2)
&\ge(1-x)(a_0-kx^3C_0)-x(1-kxr)\\
&\ge4/5-(9/5)x+2x^2-2x^3+x^4.
\end{aligned}
}{}
Its derivative is at most \(-9/5+4x<0\) on \(x\le5/12\).
Its value there is \(29309/103680>7/25\), with positive cross difference
\(6965\). If \(1-kxr<0\), the imaginary term is favorable and the real
product is at least \(\kappa[4/5-(3/2)\theta^3]>27/80>7/25\).
Also \(\alpha\ge-4703d/65536>-3d/40\), \(\beta_{a+1}\le d\), hence
\fitdisplay{
\operatorname{Cov}(C,Y)
\le-7d/25+3d^2/40\le-51d/200\le-d/4.
}{}
The case \(\alpha<0\) is included in this estimate.

\paragraph{R32.E3 Global continuous remainder with coefficient \(9\).}
Let \(W=e^{-\Gamma}Y\ge T^2F\), \(t=\sqrt d\le3/5\).
We have \(\E Y\le1\), \(Y\le q+t\), \(C\le2-q+6t/5\), and the
scaled fourth curvature is controlled by
\(H=(2-q+6t/5)[(q+1+t)^2+2]\). On \(0\le q\le2\),
\fitdisplay{
\begin{aligned}
H_0&\le32/5,\qquad
H_1=-(4/5)q^2+(22/5)q+38/5\le66/5,\\
H_2&=(7/5)q+22/5\le36/5,\qquad H_3=6/5,\quad H_4=H_5=0.
\end{aligned}
}{}
Thus \(H\le32/5+(66/5)t+(198/25)t^2\). Since \(2\Gamma\ge51t^2/16\),
\fitdisplay{
9e^{2\Gamma}-H\ge13/5-(66/5)t+(8307/400)t^2>0,
}{}
with discriminant \(-20871/500\). The resulting actual moment is
\fitdisplay{
\begin{aligned}
\E(CW^4)\le{}&e^{-4\Gamma}
 [(\E C)\mu_Y^4+4\mu_Y^3\operatorname{Cov}(C,Y)]\\
&+9e^{-2\Gamma}\operatorname{Var}Y.
\end{aligned}
}{}

\paragraph{R32.E4 The signed cubic mean factor on its actual range.}
Use \(\E C\le1-2\tau\le1-19d/16\). The bracket is at most
\(g(\mu_Y)=(1-19d/16)\mu_Y^4-d\mu_Y^3\). Its derivative is nonnegative
for \(\mu_Y\ge7/15\), because
\fitdisplay{
4(1-19d/16)\mu_Y-3d\ge4(29/48)(7/15)-1=23/180>0.
}{}
It may therefore be bounded at \(\mu_0=1-21d/40\), despite the negative
covariance. The resulting positive expression satisfies
\fitdisplay{
\mu_0^3[1-(217/80)d+(399/640)d^2]\le e^{-(2611/640)d}.
}{}
The bracket is positive, and the exponential upper uses \(d\le1/3\).
Together with \(4\Gamma\ge51d/8\), the mean/covariance exponent is
\(6691/640>52/5\). Thus
\fitdisplay{
\E[f_a\{T^2(f_3\bar f_b)\}^4]
\le F(d):=e^{-52d/5}+[(39/4)d-(243/25)d^2]e^{-51d/16}.
}{}

\paragraph{R32.E5 Exact continuous scalar payment.}
Set \(A=52/5\), \(B=51/16\), \(C=A-B=577/80\), \(b_0=39/4\),
\(b_1=243/25\). For \(0<d\le1/8\), positive Taylor through degree three
for \(e^{Bd}\), and even Taylor through degree four for \(e^{-Cd}\), imply
\fitdisplay{
\frac{e^{Bd}-e^{-Cd}}d-b_0+b_1d
\ge13/20-(1121/100)d+53d^2\ge1/20.
}{}
The unweakened quadratic coefficient is
\fitdisplay{
(B^3+C^3)/6-C^4/192=423382685439/7864320000>53,
}{}
with positive numerator difference \(6573725439\). The completed square
uses \(1121^2=1256641<1272000=10000\cdot4\cdot53\cdot(3/5)\).
Since \(e^{-Bd}>1/2\), \(F(d)\le1-d/40\) throughout this interval.
For \(1/8\le d\le1/3\),
\fitdisplay{
\begin{aligned}
F'(d)&=e^{-Bd}[P(d)-Ae^{-Cd}],\\
P(d)&=b_0-\alpha_0d+\beta_0d^2,\\
\alpha_0&=80829/1600\ge101/2,\qquad\beta_0=12393/400<31.
\end{aligned}
}{}
Put \(\Delta=Ae^{-Cd}-P\). At \(d=1/8\), \(C/8<29/32\) and the cubic
lower \(e^{-t}\ge1-t+t^2/2-t^3/6\) give
\fitdisplay{
\Delta(1/8)>(52/5)(74779/196608)-251/64=33148/983040>0.
}{}
For all \(d\) in the interval, \(C\ge7\), \(AC<609/8\),
\(\alpha_0>50\), \(\beta_0<31\), so
\fitdisplay{
\Delta'(d)\ge50-62d-\frac{609/8}{1+7d}>0.
}{}
After multiplying the positive denominator, the numerator is the concave
quadratic \(-209/8+288d-434d^2\), whose endpoint values are \(99/32\)
and \(1559/72\), both positive. Hence \(F\) decreases and
\fitdisplay{
F(d)\le F(1/8)\le1-1/320\le1-3d/320.
}{}
Combining the domains proves
\fitdisplay{
\E[f_a\{T^2(f_3\bar f_b)\}^4]\le1-3d/320
\qquad(a\ge3,b\ge4).
}{}
At \(d=0\), equality follows from \(K\one=0\); directions are independent of \(d\).

\paragraph{R32.5 Every \(L\ge3\), with actual \(h_3f_b\).}
Keep this canonical terminal for every \(b\ge4,\allowbreak L\ge3\) and every host.
Only the proof splits at \(d=1/8\). We reconstruct its finite packet,
without using an ordinary fourth-moment theorem as a same-\(f_a\) premise.

\paragraph{R32.C1 Exact low-domain packet and all path constants.}
Let \(0<d\le1/8\), \(m=b-3\), and \(G_m(-z)=A_0-zC_0\), where
\(\kappa\le A_0\le1\), \(0\le C_0\le1\), \(C_0-A_0\le x\).
The actual finite envelope is
\fitdisplay{
\begin{aligned}
F&\le e^{R-\Gamma},\qquad R=\tau u+B^3G_m(-B)u,\\
\Gamma&=3d-2\delta_2\ge35d/16,\qquad U=\tau-x^2C_0-xzA_0.
\end{aligned}
}{}
Retain \(\beta_bf_0\) and \(\tau f_1\) in the original positive
convolution for \(f_b\), use \(\beta_b\ge\tau\), and add
\(h_3=1-d+u-Bu+B^2u\). This gives \(h_3+f_b\le2-\Gamma+R\), and
nonnegative AM--GM proves the envelope without an extra original unit arm.
All coefficients are real finite polynomials. We have \(|R|<1/2\) and
\(\sigma=\lVert R\rVert_2^2<3d/20\). The actual cross
\(\E[uB^3G_m(-B)u]=-\int x^2C_0\) is nonpositive, so
\fitdisplay{
\sigma\le[(1+\theta)\theta^3+1/64]d<(55/448)d<3d/20.
}{}
The row gives \(|R|\le d+(19/16)\theta\sqrt d<1/2\) for \(d\le1/8\).
The quadratic exponential envelope with coefficient \(6/5\) yields
\fitdisplay{
\begin{aligned}
D&=D_0+(6/5)\sigma,\qquad0\le D\le47d/80,\\
D_0&=\E R=\int xA_0,\\
D_1&=\E BR=\int x^2C_0-\tau d,\\
D_2&=\E B^2R=-\int x^2A_0\le0,\\
D_3&=\E B^3R=-\int x^3C_0+\tau\delta_2\le\theta d^2.
\end{aligned}
}{}
Indeed \(D\le\theta d+(6/5)(55/448)d=31d/56<47d/80\).
The next actual means are
\fitdisplay{
\begin{aligned}
D_4&=\int x^3A_0,&D_5&=\int x^4C_0-\tau\delta_3,\\
D_6&=-\int x^4A_0,&D_7&=-\int x^5C_0+\tau\delta_4.
\end{aligned}
}{}
Each later residue class gains two powers of \(x\). Thus all \(D_i\),
\(i\ge4\), are at most \(d/14\): the first positive even term is
\(\delta_4<d/14\), the positive \(C_0\) term is at most \(\delta_5\),
and the positive \(\tau\) term at most \(\tau\delta_4\).
Put \(\epsilon=D+D_1\). Since \(D\ge0\) and \(D_1\ge-\tau d\),
\(\epsilon\ge-d^2\). The same \(R\) gives
\fitdisplay{
\begin{aligned}
\epsilon={}&\int xA_0+(1-12\tau/5)\int x^2C_0\\
&+(6/5)\lVert B^3G_m(-B)u\rVert_2^2+(-\tau+6\tau^2/5)d\\
\le{}&\delta_2+\delta_3+(6/5)(1+\theta)\delta_4\\
\le{}&(358839/524288)d<11d/16.
\end{aligned}
}{}
It also gives \(D+(D_1)_+\le11d/16\), without separately maximizing the
mean or discarding the negative cross. For the actual positive propagation,
\(Q_n=G_n+z^nU+\epsilon G_{n-2}+Dz^{n-2}\), and
\fitdisplay{
\begin{aligned}
\gamma_3&=d-D_2\ge d,\\
\gamma_4&=d+Dd-D_3\ge(243/256)d>9d/10,\\
\gamma_n&\ge\kappa d-d/14-d^3\ge(227/448)d>d/2\quad(n\ge5).
\end{aligned}
}{}
Thus \(h=1/2\) pays every required future index. The positive upper at
\(L=3\) is
\fitdisplay{
Y_3=(1+\epsilon)q-\gamma_3+BQ_2(B)u,\qquad\E Y_3\le1-97d/80.
}{}
For every \(L\ge4\), the actual positive output is bounded by
\fitdisplay{
Z_L=1+\epsilon+Q_L(B)u-(d/2)(1+q),\qquad\E Z_L\le1-13d/16.
}{}
The removed path sum is nonnegative.
Retaining the constant cross and squared mean gives
\(\operatorname{Var}Z_L\le\lVert Q_L\rVert^2d\), as in R32.2.1.

\paragraph{R32.C2 Spectral norms and receiving products.}
In the low domain,
\fitdisplay{
|U|\le d+(19/16)\theta^{3/2}<1/8+19/60=53/120<9/20.
}{}
For \(P_n=G_n+z^nU\), \(P_3\) has positive coordinates
\(1-x-x^3A_0\), \(1-\tau x+x^3C_0\). Thus
\fitdisplay{
|P_3|^2\le(1-x)^2+x(1+x^3)^2\le1,
}{}
using \(-1+x+2x^3+x^6<0\) for \(x\le1/2\). The \(G_4,\allowbreak G_5,\allowbreak G_n\)
gaps \(x/2,\allowbreak \kappa x/2,\allowbreak 7x/30\), respectively, pay the perturbations
\((9/20)x^{n/2}\). Every finite \(P_n\), \(n\ge3\), contracts, including \(n=4,\allowbreak 5\).
The added \(\epsilon+Dz\) in \(Q_3\) has norm at most \(4d/5\), because
\fitdisplay{
(11/16)^2+\theta(47/80)^2=125517/204800<16/25,
}{}
with difference \(5555/204800\). Hence \(\lVert Q_3\rVert\le1+4d/5\).
At \(n=4\), the added fee is
\([(11/16)(6/5)+(47/80)\theta]d=(2723/2560)d<(16/15)d\).
At \(n\ge5\), it is
\([11/16+(47/80)(4/15)]d=(1013/1200)d<(16/15)d\).
Thus \(\lVert Q_n\rVert\le1+(16/15)d\) for every \(n\ge4\).
The \(Q_2\) coordinates are \(1+D-\tau x+x^3C_0\), \(1+x^2A_0\),
bounded above by \(1+47/640+1/14<23/20\) and \(7/6\). Hence
\fitdisplay{
\lVert Q_2\rVert^2<529/400+637/1152<49/25,\qquad\lVert Q_2\rVert<7/5.
}{}
For the actual \(C\ge f_a\), \(\E C\le1-d\), \(C\le p+(6/5)\sqrt d\),
the receiving products give
\(\operatorname{Re}(G_aP_3)>-3/50-(9/20)(4/15)=-9/50\).
For \(n\ge4\), they are nonnegative: the \(n=4\) baseline \(57/512\)
exceeds \((9/20)\theta^2\); the \(n=5\) baseline \(1/5\) exceeds
\((9/20)\theta^{5/2}\); at \(n=6\), use its exact \(57/512\); for
\(n\ge7\), \(1451/32768\) pays the decreasing \((9/20)\theta^{n/2}\).
The tail uses the displayed lower bound, below \(1/10\).
The actual mean-product in \(\operatorname{Cov}(C,\allowbreak Q_n(B)u)\) is favorable,
and
\fitdisplay{
\operatorname{Cov}(C,Y_3)
\le9d/50+(47\theta/80)d^2+d^3\le9d/50+3d^2/8.
}{}
Negative \(\epsilon\) costs at most \(d^3\). For the long upper, the bad
product with positive \(\epsilon\) occurs only at the actual \(G_3\) index
\(L=5\), costing \((11/16)(3/50)d^2\). Negative \(\epsilon\) costs
\(d^3\), the \(D\)-power costs \((47/80)\theta d^2\), and \(-hdu\)
costs \(d^2/2\). Therefore
\fitdisplay{
\operatorname{Cov}(C,Z_L)
\le d^3+(33/800+611/2560+1/2)d^2<d^2.
}{}

\paragraph{R32.C3 Continuous remainders: \(9\) at \(L=3\), \(8\) at \(L\ge4\).}
For \(t=\sqrt d\le3/8\), \(Y_3\le q+(7/5)t+(3/8)t^2\), \(\E Y_3\le1\),
because \(\epsilon q-\gamma_3\le(11/8-1)d=3d/8\).
For \(W=e^{-\Gamma}Y_3\), the fourth remainder is at most
\(9e^{-2\Gamma}\operatorname{Var}Y_3\). Expand
\fitdisplay{
H=(2-q+6t/5)[(q+1+7t/5+3t^2/8)^2+2].
}{}
On \(0\le q\le2\),
\fitdisplay{
H_0\le32/5,\quad H_1\le27/2,\quad H_2\le21/2,\quad
H_3\le11/2,\quad H_4\le8/5,\quad H_5\le1/5.
}{}
Using \(t\le3/8\), \(H\le32/5+(27/2)t+13t^2\), and
\fitdisplay{
9e^{35t^2/8}-H\ge13/5-(27/2)t+(211/8)t^2>0,
}{}
with negative discriminant. The exact continuous coefficient calculations are
\fitdisplay{
\begin{aligned}
H_1&=-(8/5)q^2+(26/5)q+46/5,&\max H_1&=537/40<27/2,\\
H_2&=-(3/4)q^2+(43/20)q+439/50,&\max H_2&=12385/1200<21/2,\\
H_3&=669/125-(3/20)q<11/2,\\
H_4&\le1233/800<8/5,\qquad H_5=27/160<1/5.
\end{aligned}
}{}
For \(L\ge4\), \(\lVert Q_3\rVert\le1.1\), and every later norm is
at most \(17/15<6/5\). Also
\((\epsilon-d/2)q-d/2\le2(11/16-1/2)d-d/2=-d/8\), so
\(Z_L\le q+(6/5)t\), \(\E Z_L\le1\). The analogous \(H\) has
\fitdisplay{
\begin{gathered}
H_0\le32/5,\quad H_1\le66/5,\quad H_2\le216/25,\\
H_3=216/125,\qquad H_4=H_5=0.
\end{gathered}
}{}
Hence \(H\le32/5+(66/5)t+(47/5)t^2\). Retaining the positive fourth
exponential term gives
\fitdisplay{
8e^{35t^2/8}-H\ge8/5-(66/5)t+(128/5)t^2+75t^4.
}{}
For \(t\le1/6\), the quadratic decreases to \(1/9>0\).
For \(t\ge1/6\), \(75t^4\ge25t^2/12\), so the quadratic coefficient
becomes \(1661/60\), with discriminant \(-44/15\). This proves the
full low-domain coefficient \(8\).

\paragraph{R32.C4 The \(L=3\) scalar payment.}
The true mean term is at most \(e^{-73d/5}\), and the positive covariance
contribution at most \([(18/25)d+(3/2)d^2]e^{-991d/80}\).
Their sum is at most \(e^{-13d}\): after dividing the common exponential,
\fitdisplay{
e^{-49d/80}-e^{-177d/80}
\ge(8/5)d e^{-177d/80}\ge463d/400,
}{}
whereas \(18/25+3d/2\le363/400\). For negative covariance, discard that contribution.
The variance term is at most
\(9d(1+4d/5)^2e^{-35d/8}\le9de^{-111d/40}\). Thus the actual
same-\(f_a\) moment is at most \(e^{-13d}+9de^{-111d/40}\).
Put \(B=111/40\), \(C=409/40\). Positive Taylor through degree four
for \(e^{Bd}\), and even Taylor through degree six for \(e^{-Cd}\), give
\fitdisplay{
\frac{e^{Bd}-e^{-Cd}}d\ge13-(97/2)d+125d^2+700d^4.
}{}
Before weakening, the quadratic coefficient is
\(61494390880/491520000>125\), with numerator difference \(54390880\);
it equals \((B^3+C^3)/6-(C^4-B^4)/192\). The quartic coefficient is at least
\fitdisplay{
\frac{C^5(1-C/48)}{120}
>\frac{(51/5)^5(47/60)}{120}=16216186797/22500000>700.
}{}
The weaker polynomial decreases on \([0,\allowbreak 1/8]\), its derivative being at
most \(-48.5+250/8+2800/512<0\). Its endpoint \(9279/1024\) exceeds
nine by \(63/1024\). Since \(e^{-Bd}>1/2\),
\fitdisplay{
\E[f_a\{T^3(h_3f_b)\}^4]\le1-63d/2048.
}{}

\paragraph{R32.C5 Scalar payment for every \(L\ge4\).}
The mean is at most \(e^{-13d}\), and covariance at most
\(4d^2e^{-179d/16}\). Their sum is at most \(e^{-12d}\), because
\fitdisplay{
e^{-13d/16}-e^{-29d/16}\ge de^{-29d/16}\ge99d/128>4d^2.
}{}
The variance part is at most
\(8d(1+16d/15)^2e^{-35d/8}\le8de^{-269d/120}\).
Thus every actual \(L\ge4\) is bounded by \(e^{-12d}+8de^{-269d/120}\).
Set \(B=269/120\), \(C=1171/120\), \(v=C-B=451/60\in[15/2,\allowbreak 8]\).
The continuous Taylor method through degree four gives
\fitdisplay{
\frac{e^{Bd}-e^{-Cd}}d\ge12-(451/10)d+109d^2.
}{}
The unweakened quadratic coefficient is
\fitdisplay{
(B^3+C^3)/6-(C^4-B^4)/192=72+(3/2)v^2-(9/2)v-v^3/32.
}{}
It increases on \([15/2,\allowbreak 8]\), with derivative at least twelve; its left
value \(28017/256\) exceeds \(109\). The final quadratic decreases on
\([0,\allowbreak 1/8]\) to \(2581/320\), exceeding eight by \(21/320\).
Since \(e^{-Bd}>1/2\),
\fitdisplay{
\E[f_a\{T^L(h_3f_b)\}^4]\le1-21d/640
\qquad(L\ge4,\ d\le1/8).
}{}

\paragraph{R32.C6 The same direction on the high domain.}
For \(d\ge1/8\), Lemma~\ref{lem:unit-p:three-centered} gives
\(h_3f_b\le e^{-d}\) for every \(b\ge4\) in this same canonical
direction. Positivity gives
\(T^L(h_3f_b)\le e^{-d}\bar f_L\). The actual bare-reference bounds
R2-C012/A005 at \(L=3\), R2-C014/A003 at \(L=4\), and
R2-C017/A006 at \(L\ge5\), give
\fitdisplay{
\E(f_a\bar f_L^4)\le1-3d/40\qquad(a\ge3,L\ge3).
}{}
The high-domain moment is at most \(e^{-4d}(1-3d/40)\le1-3d/40\).
All low/high deficits exceed \(d/3000\). At \(d=0\), all finite tree
messages equal one because \(K\one=0\), even when \(K\ne0\).
Decorations and complete reversal use the actual positive operators and
preserve all original edges. The canonical \(h_3f_b\) direction proves
every \(L\ge3\); R32.3--R32.4 prove \(L=1,\allowbreak 2\) with \(f_3\bar f_b\).
This completes Theorem A.
